\documentclass{article}
\usepackage[utf8]{inputenc}
\usepackage{hyperref}

\title{World Models and Meta-Memory}
\author{Research Lab}
\date{\today}

\begin{document}

\maketitle

% Sections are injected by `npm run lab -- latex topics/world-models`

\section{Introduction}

The dominant path to machine reasoning scales **autoregressive language models** toward ever-larger parameter counts, accepting that inference cost and data hunger grow with capability \cite{s2_90abbc2cf38462b954ae1b772fac9532e2ccd8b0}. An alternative thesis—articulated in autonomous machine intelligence \cite{s2_775f42ed458b8c5b0f2094ea4ff5b64c557b1a34}—holds that agents should learn **predictive world models** in latent space and plan by imagination, as in early world-model RL \cite{s2_ff332c21562c87cab5891d495b7d0956f2d9228b} and modern JEPA representation learning \cite{s2_ee57e4d7a125f4ca8916284a857c3760d7d378d3}. This proposal asks whether such structured prediction can achieve **long-horizon competence with less compute and less data** than brute-scale token reasoning—and when it cannot.

\bibliographystyle{plain}
\bibliography{references}
\end{document}
